\documentclass[11pt, a4paper]{article}

\usepackage[UTF8]{ctex}
\usepackage{geometry}
\usepackage{titlesec}
\usepackage{enumitem}
\usepackage{setspace}
\usepackage{hyperref}
\usepackage{xcolor}
\usepackage{array}
\usepackage{booktabs}

\geometry{a4paper, margin=2.5cm}
\setstretch{1.25}

\hypersetup{
	colorlinks=true,
	linkcolor=black,
	urlcolor=blue
}

\title{角色卡}
\author{黄浩宇}
\date{\today}

\begin{document}
	
	\maketitle
	\newpage
	
	\tableofcontents
	\newpage
	
	\section{校长}
	
	\subsection*{左侧：内在设定}
	
	\begin{itemize}[leftmargin=2em, itemsep=0.2em]
		\item \textbf{名字 / 暂用代号}：校长
		\item \textbf{身份与背景}：学校最高决策者，善于迎合上级要求，将形式主义贯彻到学校管理的每一个环节。
		\item \textbf{性格关键词}：谄媚、专制、表演型
		\item \textbf{欲望}：获得上级的认可与表彰，让学校成为“全市领先的独一档的高中”。
		\item \textbf{伤口 / 恐惧}：害怕被上级否定、害怕学校排名下滑、害怕任何不和谐的声音出现在汇报中。
		\item \textbf{故事中的功能}：制造冲突的核心人物，将形式主义压力传导给全校，压制异议，推动荒诞规则运转。
	\end{itemize}
	
	\subsection*{右侧：外在形象}
	
	\begin{itemize}[leftmargin=2em, itemsep=0.2em]
		\item \textbf{外形与轮廓}：中等身高，体型偏瘦，站姿笔挺，习惯性地微微仰头，轮廓线条硬朗而刻板。
		\item \textbf{头发与脸}：头发梳得一丝不苟，向后贴服；脸型瘦长，颧骨略高，嘴角常年挂着一抹可随时收放的微笑。
		\item \textbf{服装与色彩}：深灰色西装套装，白色衬衫，深红色领带；主色为冷灰，辅色为暗红，风格正式而僵硬。
		\item \textbf{标志性道具 / 小动作}：讲话时双手撑在讲台上；鼓掌时率先拍三下，节奏精准；身体倾斜时始终保持头部正对前方。
	\end{itemize}
	
	\newpage
	
	\section{老师（开明老师）}
	
	\subsection*{左侧：内在设定}
	
	\begin{itemize}[leftmargin=2em, itemsep=0.2em]
		\item \textbf{名字 / 暂用代号}：老师（开明老师）
		\item \textbf{身份与背景}：持有异议的老师，是会议中唯一试图表达不同意见的人。
		\item \textbf{性格关键词}：良知、软弱、孤立
		\item \textbf{欲望}：希望学生不被占用课余时间，希望教育回归本真。
		\item \textbf{伤口 / 恐惧}：害怕被集体孤立、害怕成为众矢之的、害怕自己的声音毫无作用。
		\item \textbf{故事中的功能}：制造冲突的触发点，作为唯一异议者反衬集体机械化的荒诞，推动观众对形式主义的反思。
	\end{itemize}
	
	\subsection*{右侧：外在形象}
	
	\begin{itemize}[leftmargin=2em, itemsep=0.2em]
		\item \textbf{外形与轮廓}：中等身材，略微驼背，站姿与周围人格格不入——当所有人倾斜时，他试图坐直。
		\item \textbf{头发与脸}：头发略显凌乱，没有刻意打理；面容疲惫，眼神中带有犹豫和不安。
		\item \textbf{服装与色彩}：浅蓝色衬衫，袖口微皱，没有领带；主色为柔和的蓝，辅色为灰白，风格朴素而日常。
		\item \textbf{标志性道具 / 小动作}：欲言又止时嘴唇微动；被注视时慢慢低下头；试图坐直身体却显得僵硬。
	\end{itemize}
	
	\newpage
	
	\section{教育部部长}
	
	\subsection*{左侧：内在设定}
	
	\begin{itemize}[leftmargin=2em, itemsep=0.2em]
		\item \textbf{名字 / 暂用代号}：教育部部长
		\item \textbf{身份与背景}：上级视察领导，重视汇报成果与表面评价，满足于被安排的“标准答案”。
		\item \textbf{性格关键词}：官僚、虚荣、易满足
		\item \textbf{欲望}：看到完美的汇报成果，获得“教育先进市”的政绩。
		\item \textbf{伤口 / 恐惧}：害怕视察中出现纰漏、害怕汇报数据不好看、害怕下级不给面子。
		\item \textbf{故事中的功能}：推动剧情进入视察环节，以“被安排的满意”强化形式主义闭环，制造讽刺效果。
	\end{itemize}
	
	\subsection*{右侧：外在形象}
	
	\begin{itemize}[leftmargin=2em, itemsep=0.2em]
		\item \textbf{外形与轮廓}：身材高大，体态微胖，站姿自信而放松，习惯性地将双手背在身后或叉腰，轮廓圆润而有压迫感。
		\item \textbf{头发与脸}：头发稀疏，梳得油亮贴服；圆脸，笑容来得快收得也快，眼神中带着审视与满意的切换。
		\item \textbf{服装与色彩}：深藏青色中山装式外套，金属扣子，内搭白衬衫；主色为深蓝，辅色为银白，风格庄重而保守。
		\item \textbf{标志性道具 / 小动作}：提问时微微前倾；满意时拍对方肩膀；笑容凝固时缓缓转头。
	\end{itemize}
	
	\newpage
	
	\section{秘书}
	
	\subsection*{左侧：内在设定}
	
	\begin{itemize}[leftmargin=2em, itemsep=0.2em]
		\item \textbf{名字 / 暂用代号}：秘书
		\item \textbf{身份与背景}：部长的随行人员，负责执行部长的指示，包括拍照记录视察“成果”。
		\item \textbf{性格关键词}：顺从、沉默、工具化
		\item \textbf{欲望}：不出差错地完成部长交代的任务。
		\item \textbf{伤口 / 恐惧}：害怕拍照时机不对、害怕被部长责备。
		\item \textbf{故事中的功能}：陪伴型角色，以相机快门声完成“定格”的仪式感，强化形式主义被记录、被传播的荒诞。
	\end{itemize}
	
	\subsection*{右侧：外在形象}
	
	\begin{itemize}[leftmargin=2em, itemsep=0.2em]
		\item \textbf{外形与轮廓}：中等身高，体型匀称，站姿恭敬而克制，始终站在部长侧后方半步的位置，轮廓低调。
		\item \textbf{头发与脸}：短发整齐，面容平静，表情始终保持中性，眼神跟随部长移动。
		\item \textbf{服装与色彩}：黑色西装，白衬衫，黑色领带；主色为纯黑，辅色为白，风格极简而正式。
		\item \textbf{标志性道具 / 小动作}：手持相机，随时准备拍摄；部长示意时立即上前；拍照时动作迅速而精准。
	\end{itemize}
	
	\newpage
	
	\section{学生A}
	
	\subsection*{左侧：内在设定}
	
	\begin{itemize}[leftmargin=2em, itemsep=0.2em]
		\item \textbf{名字 / 暂用代号}：学生A
		\item \textbf{身份与背景}：能够熟练说出标准答案的学生，被安排好在视察路线上“偶遇”部长。
		\item \textbf{性格关键词}：训练有素、表演型、无自觉
		\item \textbf{欲望}：完成被安排的回答任务，获得上级的认可。
		\item \textbf{伤口 / 恐惧}：害怕说错答案、害怕让校长失望。
		\item \textbf{故事中的功能}：制造冲突与讽刺的核心角色，以“抢答标准答案”揭示形式主义下学生被规训的荒诞。
	\end{itemize}
	
	\subsection*{右侧：外在形象}
	
	\begin{itemize}[leftmargin=2em, itemsep=0.2em]
		\item \textbf{外形与轮廓}：青少年身形，偏瘦，站姿笔直而僵硬，动作幅度大而精准，像被训练过的表演者。
		\item \textbf{头发与脸}：短发整齐，面容清秀，脸上挂着标准而固定的微笑，眼神明亮却缺乏变化。
		\item \textbf{服装与色彩}：蓝白相间的校服，拉链拉到固定高度，白色运动鞋；主色为蓝白，风格统一而刻板。
		\item \textbf{标志性道具 / 小动作}：抢先一步向前跨出；声音洪亮地背诵答案；保持微笑一动不动；被拍肩膀时笑容不变。
	\end{itemize}
	
	\newpage
	
	\section{群众领导、群众老师}
	
	\subsection*{左侧：内在设定}
	
	\begin{itemize}[leftmargin=2em, itemsep=0.2em]
		\item \textbf{名字 / 暂用代号}：群众领导、群众老师
		\item \textbf{身份与背景}：参加会议的人员，构成会议现场的“集体背景”，以整齐划一的动作配合校长。
		\item \textbf{性格关键词}：从众、机械、无个性
		\item \textbf{欲望}：跟随校长的节奏，不出错、不突出。
		\item \textbf{伤口 / 恐惧}：害怕成为异类、害怕被集体排斥。
		\item \textbf{故事中的功能}：制造集体压力的背景力量，以机械化的同步动作反衬开明老师的孤立，强化荒诞氛围。
	\end{itemize}
	
	\subsection*{右侧：外在形象}
	
	\begin{itemize}[leftmargin=2em, itemsep=0.2em]
		\item \textbf{外形与轮廓}：群体形象，身高体型各异但姿态高度统一，坐姿整齐，身体可同步倾斜四十五度。
		\item \textbf{头发与脸}：发型各异但梳理整齐；面部表情统一为严肃而空洞，笑容可同步出现与收住。
		\item \textbf{服装与色彩}：深色正装为主，颜色统一在灰、黑、深蓝范围内；风格正式而缺乏变化。
		\item \textbf{标志性道具 / 小动作}：身体统一向一侧倾斜四十五度；整齐拍手三下；同步转头；整齐笑声“哈哈哈”。
	\end{itemize}
	
	\newpage
	
	\section{群众学生}
	
	\subsection*{左侧：内在设定}
	
	\begin{itemize}[leftmargin=2em, itemsep=0.2em]
		\item \textbf{名字 / 暂用代号}：群众学生
		\item \textbf{身份与背景}：参加周六兴趣班的学生群体，在报时器的指令下统一行动。
		\item \textbf{性格关键词}：服从、机械、被规训
		\item \textbf{欲望}：完成被安排的流程，跟随指令行动。
		\item \textbf{伤口 / 恐惧}：害怕不整齐、害怕脱离集体节奏。
		\item \textbf{故事中的功能}：制造荒诞氛围的群体角色，以“兴趣班”的机械流程讽刺形式主义对学生的规训。
	\end{itemize}
	
	\subsection*{右侧：外在形象}
	
	\begin{itemize}[leftmargin=2em, itemsep=0.2em]
		\item \textbf{外形与轮廓}：青少年群体，坐姿与站姿高度统一，动作如机器人般精准一致。
		\item \textbf{头发与脸}：发型统一符合校规；面无表情，视线统一朝向讲台。
		\item \textbf{服装与色彩}：统一校服，蓝白配色，整齐划一。
		\item \textbf{标志性道具 / 小动作}：双手整齐放在桌面上；听到“滴”声同时低头通电；同时起立；迈机器人步伐离开。
	\end{itemize}
	
	\newpage
	
	\section*{注}
	
	以上角色卡根据剧本中角色的台词、动作描写、场景功能及象征意义生成。校长、老师、部长、秘书、学生A为具名个体角色；群众领导/老师、群众学生为群体角色，其外在形象以“统一性”为核心特征。
	
\end{document}